% Meta: 100–200 palabras
\section{Justificación}

Anticipar el rendimiento por parcela, antes o al cierre del ciclo, cambia la calidad de la decisión. Un productor, un técnico o un intermediario financiero deja de depender solo del recuerdo de la zona y puede contrastar un estimado con el clima de la temporada y con el estado del cultivo visto desde satélite. En cebada de temporal, en el altiplano de Hidalgo, Puebla y Tlaxcala, esa información es escasa a escala de parcela y llega tarde si solo existe al momento de la cosecha.

El reto existe porque los datos públicos de percepción remota y de clima ya cubren la región, pero no están convertidos en una herramienta de consulta. Justifica el trabajo traducirlos a toneladas por hectárea, con un error medido y con límites explícitos, para que el número sirva de apoyo y no de promesa. Un estimado honesto, aunque imperfecto, es más útil para el crédito y la asistencia técnica que un modelo opaco que no dice dónde falla.
